\documentclass{beamer}

\usepackage[russian]{babel}
\usepackage[utf8]{inputenc}
\usepackage{cmap}
\usepackage{graphicx}
\usepackage{xspace}
\usepackage{psfrag}
 \usepackage{hyperref}
\usepackage[normalem]{ulem}

\beamertemplatenavigationsymbolsempty
\setbeamertemplate{footline}[frame number]
\usecolortheme{seahorse}

%%Исправляем греческие буквы в формулах, они должны быть прямыми, а не курсивными
\DeclareSymbolFont{greek}{U}{eur}{m}{n}                                                        
\SetSymbolFont{greek}{bold}{U}{eur}{b}{n}                                                      
\DeclareSymbolFontAlphabet{\gr}{greek}                                                         
\DeclareMathSymbol{\alpha}{\mathord}{greek}{"0B}                                               
\DeclareMathSymbol{\beta}{\mathord}{greek}{"0C}                                                
\DeclareMathSymbol{\gamma}{\mathord}{greek}{"0D}                                               
\DeclareMathSymbol{\delta}{\mathord}{greek}{"0E}                                               
\DeclareMathSymbol{\epsilon}{\mathord}{greek}{"0F}                                             
\DeclareMathSymbol{\zeta}{\mathord}{greek}{"10}                                                
\DeclareMathSymbol{\eta}{\mathord}{greek}{"11}                                                 
\DeclareMathSymbol{\theta}{\mathord}{greek}{"12}                                               
\DeclareMathSymbol{\iota}{\mathord}{greek}{"13}                                                
\DeclareMathSymbol{\kappa}{\mathord}{greek}{"14}                                               
\DeclareMathSymbol{\lambda}{\mathord}{greek}{"15}                                              
\DeclareMathSymbol{\mu}{\mathord}{greek}{"16}                                                  
\DeclareMathSymbol{\nu}{\mathord}{greek}{"17}                                                  
\DeclareMathSymbol{\xi}{\mathord}{greek}{"18}                                                  
\DeclareMathSymbol{\pi}{\mathord}{greek}{"19}                                                  
\DeclareMathSymbol{\rho}{\mathord}{greek}{"1A}                                                 
\DeclareMathSymbol{\sigma}{\mathord}{greek}{"1B}                                               
\DeclareMathSymbol{\tau}{\mathord}{greek}{"1C}                                                 
\DeclareMathSymbol{\upsilon}{\mathord}{greek}{"1D}                                             
\DeclareMathSymbol{\phi}{\mathord}{greek}{"1E}                                                 
\DeclareMathSymbol{\chi}{\mathord}{greek}{"1F}                                                 
\DeclareMathSymbol{\psi}{\mathord}{greek}{"20}                                                 
\DeclareMathSymbol{\omega}{\mathord}{greek}{"21}                                               
\DeclareMathSymbol{\varepsilon}{\mathord}{greek}{"22}                                          
\DeclareMathSymbol{\vartheta}{\mathord}{greek}{"23}                                            
\DeclareMathSymbol{\varpi}{\mathord}{greek}{"24}                                               
\let\varrho\rho                                                                                
\let\varsigma\sigma                                                                            
\DeclareMathSymbol{\varphi}{\mathord}{greek}{"27}

\title{Разработка интеллектуального ассистента для~планирования путешествий}
\author{Мицкевич Валерия Олеговна}
\date{\today}

\begin{document}
\maketitle

\begin{frame}
\frametitle{Актуальность работы}
\framesubtitle{Проблематика, объект и предмет  исследования}

Интеллектуальные ассистенты значительно упрощают процесс планирования поездок, предоставляя актуальную информацию о~погоде,
транспорте, местах отдыха и отелях.
Они могут автоматически составлять маршруты, учитывая пожелания пользователя, сформулированные в~его запросе.
Это экономит время и силы, а~также позволяет избежать стресса при~организации путешествия.

{\bf Объект исследования:} механизм tools calling для~связи языковой модели и интернет-сервисов.

{\bf Предмет исследования: } выполнение запросов к~интернет-сервисам при~помощи языковых моделей.
\end{frame}

\begin{frame}
\frametitle{Цель и задачи}

{\bf Цель:} разработать интеллектуальный ассистент в~предметной области организации путешествий для~обработки запросов пользователя, сформулированных на~естественном языке.

{\bf Задачи:}

\begin{enumerate}
\item {Организация взаимодействия с~интернет-сервисами для~получения информации о~погоде и планирования поездки.}
\item {Организация взаимодействия с~языковой моделью с~использованием технологии tools calling.}
\item {Тестирование и оценка полученных результатов.}
\end{enumerate}
\end{frame}

\begin{frame}
\frametitle{Обработка запроса}
\framesubtitle{последовательность обмена информацией для~получения ответа}
\begin{figure}[h]
\includegraphics[width=9cm,height=5cm,keepaspectratio]{tools-seq.pdf}
\end{figure}
\end{frame}

\begin{frame}
\frametitle{Компоненты системы}
\framesubtitle{Структура взаимодействия компонентов системы}
\begin{figure}[h]
\includegraphics[width=9cm,height=5cm,keepaspectratio]{service-comp.pdf}
\end{figure}
\end{frame}

\begin{frame}
  \frametitle{Заключение}
  \framesubtitle{Текущие результаты работы}

  \begin{enumerate}
  \item {
Разработан интеллектуальный ассистент для~планирования поездок с~обработкой запросов на~естественном языке. 
  }
  \item {
Организовано взаимодействие с~сервисами YandexGPT, Яндекс.Расписания, Open-Meteo.
  }
  \item {
Работа ассистента протестирована на~проверочных примерах с~получением планируемого результата.
    }
    \end{enumerate}
\end{frame}

\begin{frame}
  \begin{center}
{\Large Спасибо за~внимание!}
    \end{center}
  \end{frame}



\end{document}
